\clearpage
\begingroup
\fancyhead[L]{\small NACHLASS UND BRIEFWECHSEL}
\renewcommand{\headrulewidth}{0pt}
\SourceLeaf{605}
\par\medskip
\Needspace{8\baselineskip}
\pdfbookmark[2]{[7]}{pallas-IX-fragment7}
\begin{center}[7.]\end{center}
\begingroup
\small
\par\noindent\hspace*{\parindent}[Die Berechnung der Grössen \(M,N,\varphi\) nach den Gleichungen 3 der
Notiz [5] muss durch Annäherung erfolgen; man kann indessen die Wurzeln der
Gleichung \(\rho\rho=0\) auch auf folgende Weise bestimmen:]
\par\noindent\hspace*{\parindent}Auflösung der Gleichung
\[
a+b\cos2u=\cos(u+D)
\]
\noindent[wo \(u\) für \(E'\) geschrieben ist und die Beziehungen der Constanten \(a,b,D\)
zu denen der Gleichung 1 der Notiz [5] leicht ersichtlich sind]:
\par\noindent\hspace*{\parindent}Man setze
\[
u=x+y,
\]
\noindent also
\[
\begin{aligned}
a+b\cos2x\cos2y-\cos(D+x)\cos y&\\
{}+\sin y(\sin(D+x)-2b\sin2x\cos y)&=0.
\end{aligned}
\]
\par\noindent\hspace*{\parindent}Setzen wir also, was erlaubt ist,
\[
\text{1)}\qquad 2b\sin2x\cos y=\sin(D+x),
\]
\noindent so wird
\[
\text{2)}\qquad b\cos2x\cos2y-\cos(D+x)\cos y+a=0.
\]
\par\noindent\hspace*{\parindent}Die Substitution des Werths von \(\cos y\) aus 1) in 2) gibt
\[
4b\sin2x^2\mathord{.}(a-b\cos2x)=\cos2D-\cos2x,
\]
\noindent also, wenn man \(\cos2x=t\) setzt,
\[
\text{3)}\qquad 4b(1-t)(1+t)(a-bt)=\cos2D-t
\]
\noindent oder entwickelt
\[
\text{4)}\qquad 4bbt^3-4abtt+(1-4bb)t+4ab-\cos2D=0.
\]
\par\noindent\hspace*{\parindent}Die Form 3) ist zur indirecten Auflösung am bequemsten.
\par\noindent\hspace*{\parindent}Nachdem man also \(x\) gefunden hat, kann man \(y\) entweder aus 1) durch
die Formel
\[
\cos y=\frac{\sin(D+x)}{2b\sin2x}
\]
\SourceLeaf{606}
\noindent oder aus 2) durch die Formel
\[
\cos y=\frac{(-a+b\cos2x)\sin2x}{\sin(D-x)}
\]
\noindent bestimmen. Für \(x\) kann man gewiss sein, allemal einen reellen Werth zu
erhalten. Fände man für \(\cos y\) einen Werth grösser als 1, so ist der imaginäre
Werth von \(y\):
\[
y=i\log\mathrm{hyp}\{\cos y+\BareRoot{(\cos y^2-1)}\}.
\]
\noindent Bedeutet \(u\) eine excentrische Anomalie und \(e\) die Excentricität, so ist
die entsprechende mittlere Anomalie
\[
\begin{aligned}
x-e\sin x\cos y+i\{&\log\mathrm{hyp}(\cos y+\BareRoot{(\cos y^2-1)})\\
&-e\cos x\BareRoot{(\cos y^2-1)}\}.
\end{aligned}
\]
\par\noindent\hspace*{\parindent}Sind
\[
a\cos(nM+A),\qquad a'\cos((n+1)M+A')
\]
\noindent zwei auf einander folgende Glieder einer Reihe, die eine unendlich gross
werdende Function vorstellt, wenn \(a+b\cos2u=\cos(u+D)\) würde, so ist
\[
\begin{gathered}
\mathrm{Limes}\,(A-A')=x-e\sin x\cos y\\
\mathrm{Limes}\,\frac a{a'}=\frac{\cos y+\BareRoot{(\cos y^2-1)}}{e\cos x\BareRoot{(\cos y^2-1)}},
\end{gathered}
\]
\noindent wo \(c\) die Basis der hyperbolischen Logarithmen bedeutet. Zur bequemern
Rechnung kann man noch \(\cos y=\frac1{\sin\zeta}\) setzen, wodurch
\[
\mathrm{Limes}\,\frac a{a'}=\frac{\mathrm{cotang}\,\frac12\zeta}{c^{e\cos x\,\mathrm{cotang}\,\zeta}}
\]
\noindent wird.
\par\noindent\hspace*{\parindent}[Man wird also, übereinstimmend mit den vorigen Notizen, für den Factor
\(\lambda=\left(\frac1{\sin\zeta'}-\cos(M'-\varphi')\right)^{\frac32}\) die Werthe annehmen,
die der Wurzel der Gleichung \(\lambda=0\) entsprechen:
\[
\varphi'=x-\frac{e\sin x}{\sin\zeta}=\mathrm{Limes}\,(A-A')
\]
\[
\begin{gathered}
\frac1{\sin\zeta'}=\cos\left(i\log\mathrm{hyp}\,\mathrm{Limes}\,\frac a{a'}\right)\\
\text{oder}\qquad \mathrm{cotang}\,\tfrac12\zeta'=\mathrm{Limes}\,\frac a{a'}\mathpunct{\cdot}\Biggr]
\end{gathered}
\]
\par\medskip\begin{center}\rule{4em}{.4pt}\end{center}
\endgroup

\clearpage
\endgroup
